\section{Related work}
\label{sec:related}

The construction we transform is that of Jones, Sato, Wada and Wiens
\cite{jsww1976}, itself built on the solution of Hilbert's tenth problem
\cite{matiyasevich1970} and, for their Section~3, on the reduction of
Matijasevi\v{c} and Robinson \cite{mr1975}. The frontier at the low-variable end
is due to Matiyasevich \cite{matiyasevich1981}, whose generator in $10$
variables remains the fewest known and is, unlike most figures in this range,
genuinely exhibited.

Universal Diophantine pairs $(\nu,\delta)$---bounds valid for every Diophantine
set---are tabulated by Jones \cite{jones1982}, who also supplied the complete
proof of the nine-unknowns theorem that Matiyasevich had announced. These pairs
are frequently quoted as though they were statements about the primes; they are
not, and \S\ref{sec:correc} discusses the confusion this causes. Over $\Z$ the
corresponding record is due to Sun \cite{sun2021}. Bayer and David
\cite{bayerdavid2025} formalize complexity bounds for Diophantine equations and
note that some universal pairs in circulation lack published proofs.

Closest to the present paper is the work of P\k{a}k and Kaliszyk, who formalized
in Mizar both the $(26,25)$ system of \cite{jsww1976} \cite{pak2022jsww} and
Matiyasevich's polynomial in ten unknowns \cite{pak2022ten}. Their work
formalizes the two \emph{endpoints} of the known frontier; ours exhibits and
verifies points \emph{between} them, and in addition audits the announcements
that connect them. The two efforts are complementary, and their figure for the
degree of the $12$-variable polynomial ($13697$) is one of the four numbers our
implementation of \cite{mr1975} reproduces.

\section{Conclusion}
\label{sec:concl}

We have exhibited two generators of the set of primes, at $(44,5)$ and
$(21,25)$, together with machine-checked proofs that each is equisatisfiable
with the system published in \cite{jsww1976}. We did not find a prime generator of degree $5$ written out anywhere in the
literature we surveyed; the second improves by five variables on the only
previously exhibited generator at that degree. We have
also supplied the fourteen substitution steps that \cite{jsww1976} states
without proof for their quotient-method system, and corrected three items in the
secondary literature.

We have not improved on any record. In the number of variables the
$(10,>6000)$ of \cite{matiyasevich1981} has fewer than anything our method
reaches, and going below it would require the nine-unknowns machinery, which our
approach---transcribe a published system, eliminate, flatten, and measure---does
not reproduce. In the degree, $5$ is the floor for generators obtained by the
construction of Proposition~\ref{prop:gen}; whether a prime generator of degree
$3$ or $4$ of some other shape exists we do not know. So the only question we
can address there is the variable count, and we do not know whether $44$ can be
lowered.

What the exercise does suggest is that this part of the literature is less
thoroughly worked than its age implies. Figures announced and never written;
a method cited for a figure it does not yield; two conflations of units in wide
circulation; and fourteen substitution steps whose soundness over $\N$ had not
been checked. None of that required new mathematics to find---only the decision
to treat every announced pair as something to be exhibited rather than quoted.

\subsection*{Availability}

The systems, the optimizer, the Lean development and the tests that check each
formal statement against the published equations are available at
\url{https://github.com/Javitax47/Proyecto-Diophantus}. The \LaTeX{} of the
system in \S\ref{ssec:flat-system} is generated from the same file the Lean
kernel checks.
